\section{Structuring lemmas for reuse}\label{sec:13-working-efficiently:05-structuring-lemmas:1}

The single biggest efficiency gain, greater than any tactic choice, is to \textbf{prove the general fact once, as its own named lemma, as soon as an argument would otherwise be repeated.} \texttt{left\_\allowbreak{}inverse\_\allowbreak{}unique} from Chapter 8 is the running example. Theorem 3 (\texttt{inv\_\allowbreak{}op}) and \texttt{neg\_\allowbreak{}one\_\allowbreak{}mul} and \texttt{neg\_\allowbreak{}mul} from Chapter 10 all reduce to it, instead of re-deriving ``uniqueness of inverses'' inline. Signs that a lemma should be factored out are as follows.

\begin{itemize}
\tightlist
\item
  A \href{https://lean-lang.org/doc/reference/latest/Tactic-Proofs/Tactic-Reference/}{\texttt{rw}} chain already written for a different but structurally identical goal is about to be repeated. In that case, the shared shape should instead be stated as its own \texttt{theorem}/\texttt{have}, then applied via \texttt{apply}/\texttt{exact} in both places.
\item
  A sub-goal deep in a proof would itself be a reasonable, independently statable mathematical fact (for example, ``an element that equals its own double is zero,'' buried inside \texttt{mul\_\allowbreak{}zero} of Chapter 10). Naming it is both more efficient \emph{and} more readable, since the outer proof then reads as a short chain of named facts instead of one long undifferentiated block.
\end{itemize}

This is the same judgment call made writing ordinary code, extract a helper when, and only when, real duplication or a genuinely separate sub-claim is present, not ahead of time ``just in case.''

\textbf{Key points.} \texttt{exact?}/\texttt{apply?} search for a closing term but do not always find the shortest one. \texttt{decide}/\texttt{omega}/\texttt{norm\_\allowbreak{}num} replace a hand proof exactly on their decidable fragment, never on a goal about a generic, unspecified structure. \texttt{simp} trades traceability for speed, so this book reaches for it only when a genuine technical obstacle makes the explicit alternative not worth the detour. And a repeated \texttt{rw} chain or an independently-statable sub-goal is a lemma waiting to be named.

\textbf{Socratic questions.}

\begin{enumerate}
\def\labelenumi{\arabic{enumi}.}
\tightlist
\item
  \emph{\texttt{exact?} sometimes returns a working but needlessly roundabout term instead of the shortest one (the worked example of Chapter 13, Section 1 itself). Does that make it a poor tool to reach for?} No.~It still guarantees a \emph{correct} closing term, found automatically. Simplifying what it returns by hand afterward is a small, separate step, cheaper than deriving the whole thing from nothing.
\item
  \emph{\texttt{decide} and a hand-written proof can close the exact same goal on a finite carrier like \texttt{Fin 3}. \texttt{fin3Group}/\texttt{fin3Ring} from Chapter 9 already used \texttt{decide} this way. If it was already in use back then, what does this chapter actually add?} Chapter 9 used \texttt{decide} on one specific example, without stopping to say \emph{why} it applied there. This chapter turns that one-off use into a general principle, recognizing which regime a goal is in, a finite, concretely enumerable carrier where \texttt{decide} applies, versus a general, unspecified structure (an arbitrary \texttt{Group G} or \texttt{Ring R}, as in Chapters 8 and 10) where no decision procedure applies and a hand-built proof is the only option. Knowing \emph{when} to reach for \texttt{decide} is the skill; using it once, correctly, is not yet that skill.
\item
  \emph{A repeated \texttt{rw} chain is a sign a lemma should be factored out. Is the reverse also true, should every proof be split into the smallest possible named lemmas, just in case one gets reused later?} No. Speculative extraction with no present duplication is the same mistake in the opposite direction, three similar lines are better than a premature abstraction. The signal to extract is an actual repeated shape or an actual independently-statable sub-claim, not the mere possibility of future reuse.
\end{enumerate}

